\documentclass[11pt]{article}
\usepackage[T1]{fontenc}
\usepackage[utf8]{inputenc}
\usepackage[a4paper,margin=27mm]{geometry}
\usepackage{lmodern,amsmath,amssymb,amsthm,microtype,hyperref}
\hypersetup{hidelinks}
\setlength{\emergencystretch}{3em}
\newtheorem{theorem}{Theorem}[section]
\newcommand{\workbenchurl}{https://github.com/KokunoYumeto/zeta-and-l-functions-workbench}
\newcommand{\recordurl}{https://github.com/KokunoYumeto/mathematics-workbenches}
\providecommand{\tightlist}{\setlength{\itemsep}{0pt}\setlength{\parskip}{0pt}}
\title{Equivariant theta transport after CocoSgt:\\the cubic Riemann hypothesis}
\author{}
\date{24 November 2026}
\expandafter\def\expandafter\UrlBreaks\expandafter{\UrlBreaks\do\a\do\b\do\c\do\d\do\e\do\f\do\g\do\h\do\i\do\j\do\k\do\l\do\m\do\n\do\o\do\p\do\q\do\r\do\s\do\t\do\u\do\v\do\w\do\x\do\y\do\z\do\A\do\B\do\C\do\D\do\E\do\F\do\G\do\H\do\I\do\J\do\K\do\L\do\M\do\N\do\O\do\P\do\Q\do\R\do\S\do\T\do\U\do\V\do\W\do\X\do\Y\do\Z\do\0\do\1\do\2\do\3\do\4\do\5\do\6\do\7\do\8\do\9}
\begin{document}
\maketitle
\begin{abstract}
CocoSgt's pinned formalization of the theta framework supplies a selected-lattice gain that is a uniform image inclusion. We identify the exact local quantifiers and prove that a compatible averaging projector preserves the gain exponent. For full lattices of finite index, a coherent direct sum of transported copies preserves both the gain and the normalized lattice cost, and its invariant projection recovers the original system. We compute its projected prime-power trace and prove a scalar-commutant obstruction for a natural orbit of distinct highest weights. The local alternating weights force odd cyclic shifts to exchange two disjoint source windows; dual lattice transport changes the Hecke edge type from $r$ to $k-r$. Adjoining the source's CM quadratic field does not remove that local shift, and the literal size assumptions force that quadratic field to vary. These conclusions concern the specified constructions and conditional source interfaces, and do not prove an obstruction to every equivariant theta method or a new prime-cancellation estimate. A separate statement-only reconstruction verifies these scoped claims while retaining the source realization hypotheses. The mathematics and text were produced by GPT-6 Astra in ChatGPT Pro mode; no human has checked this text.
\end{abstract}

\section{Main statements and exact scope}
Fix the source commit
\path|7cbf7a0c32fca1badf05931879a37c3f8a2b03b4|
in the repository \cite{source}. Its local gain interface uses a discrete
valuation ring $O$, fraction field $E$ of characteristic zero, uniformizer $\pi$, an original
lattice $M$, a selected sublattice $N=\mathrm{NP}$, a linear theta-chain
map $f$, and a target lattice $B=\mathrm{Int}_m$. The source-specific
gain hypothesis is an inclusion $f(N)\subseteq\pi^a B$, with the
same integer $a$ for every vector in $N$; its printed realization
hypotheses remain assumptions throughout.

All vector spaces in the independent linear-algebra statements below
are finite-dimensional over $E$. A full lattice is a finitely generated
$O$-submodule spanning its ambient vector space. Whenever a quotient
length is divided by rank, $M_0$ is a full lattice of positive rank,
$N\subseteq M_0$ is full, and $M_0/N$ has finite length. These are
hypotheses of the normalized-index statements; the source's module
and rank conventions are stated separately in Section~2.

\begin{theorem}
The following statements hold.
\begin{enumerate}
\item If a finite group $H$ acts $E$-linearly, preserves $N,B$, and
$f$ intertwines its actions, then
\[
f(e_HN)\subseteq\pi^a e_HB,\qquad
e_H=|H|^{-1}\sum_{h\in H}h.
\]
There is no exponent loss. The analogous inference for a
length bound divided by rank is false in general.
\item Given isomorphisms $T_{h,j}:V_j\to V_{h,j}$, the maps
$U_a|_{V_{h,j}}=T_{ah,j}T_{h,j}^{-1}$ define a coherent action on
the direct sum of transported copies. Transported chain maps commute
with this action. The projected lattice system is isomorphic to the
original lattice system, including its normalized finite index under the full-lattice and finite-length hypotheses above.
For a finite abelian group $A$, $H\subset A$, and
$\mathcal F_p=\chi(p)L_{a_p}\otimes T_p$, one has
\[
\operatorname{tr}(e_H\mathcal F_p^\nu)
=|A/H|\,\mathbf{1}_{a_p^\nu\in H}\,
\chi(p)^\nu\operatorname{tr}(T_p^\nu).
\]
\item Over an algebraically closed field of characteristic zero,
the direct sum of the distinct irreducible highest-weight modules
in a transitive permutation orbit has scalar joint commutant for
the full coefficient group $\prod_\tau\mathrm{GL}_k$ and the
natural permutation action. Thus that joint commutant contains
no proper nonzero idempotent.
\item In the source's even cyclic local coordinates, strict
transport of the effective weight requires
$\lambda^{(c)}_{s+c}=J^{c\bmod2}\lambda_s$, where
$J\lambda=-\operatorname{reverse}(\lambda)$. Under the full literal
family windows, an odd shift of a top position exchanges
$a_s,b_s$ with $b_s<a_s$ and leaves those windows.
Dualizing a lower type-$r$ edge yields an upper type-$r$ edge,
or a lower type-$(k-r)$ edge after central rescaling.
\item At a source inert-class prime, the local step is
$\operatorname{Frob}_p c=((1,1),h)$ in the source's Galois
coordinates, under its class/Frobenius comparison interface.
It already fixes the quadratic CM base $K_d$.
Moreover the literal rank and size assumptions imply
\[
d>16x^{\,4+799\epsilon/50}.
\]
Thus adjoining $K_d$ does not remove an offending local step
that lies in the projector subgroup, and $d$ cannot be fixed
along an unbounded literal family.
\end{enumerate}
\end{theorem}

The last statement does not assert the existence of such an
offending prime below $2x$. The highest-weight and edge statements
concern the specified natural transports. More general theta-chain
intertwiners are not ruled out. The finite general-linear operators
in the source are not identified here with a field-deck action.

\section*{Proof overview}
The source's valuation predicate is literal membership in a scalar multiple of a lattice, with the same exponent for every input in its selected sublattice. Consequently an intertwining averaging projector preserves that inclusion. An induced direct sum makes every system equivariant by transport, but a change of coordinates identifies it with a regular representation tensored with the original system. This identifies both its projected lattice and its Frobenius trace, rather than merely asserting that arithmetic compatibility is absent. For the natural orbit of nonisomorphic highest weights, the full coefficient group and the permutation action have scalar joint commutant. In the local source coordinates, an odd cyclic shift additionally requires dualizing the highest weight; the source windows are not closed under that operation, and lattice duality reverses the Hecke edge type. Finally, the proposed enlargement by the source's CM quadratic field neither changes the relevant local step nor allows that quadratic field to stay fixed at unbounded scales.

\section{The exact source gain and its quantifiers}\label{the-exact-source-gain-and-its-quantifiers}

Write \(O\) for the source's local discrete valuation ring, \(E\) for its characteristic-zero fraction field, \(p\) for the rational residual characteristic, and \(e_p=v_O(p)\). Fix a uniformizer \(\pi\). At one admissible word and place, put

\[
 M=G.M,\qquad M_{\rm ram}=M\cap V_{\rm ram},\qquad
 f=\operatorname{chainK}_m,\qquad B=\operatorname{Int}_m.
\]

The source's \path|vGeLat| and \path|LatInst.vKge| mean precisely

\[
 \operatorname{vKge}(F,g)
 \quad\Longleftrightarrow\quad
 f(F)\in\pi^{\lceil e_p g\rceil}B.\tag{2.1}
\]

Although \path|vGeLat| quantifies an irreducible element separately for each assertion, this makes no difference: two uniformizers differ by a unit, and multiplication by a unit preserves an \(O\)-lattice, also after multiplication by a negative integral power of a uniformizer. The source proves this conversion explicitly in \path|TopChoiceBounds.bg_vGe_iff|.

Let \(VV\) be the selected set of good orbit representatives, \(y_{V'}\) the ordered product of the designated operators for \(V'\subseteq VV\), and \(S_v\) its source cost. The source lattice is

\[
 \begin{split}
 I&=\operatorname{span}_O\left\{
 y_{V'}\!\left(p^{\sum_{v\in VV\setminus V'}\max(0,\lceil S_v\rceil)}F\right):
 V'\subseteq VV, F\in M\right\},\\
 N&=I\cap p^{\max(0,\lceil\tau_P\rceil)}M\cap V_{\rm ram}.
 \end{split}\tag{2.2}
\]

Here \(I\) is \path|IMP| and \(N\) is \path|NP|. The finite product is the source's specified product; no unproved commutativity is being imposed. Its prime gain package \path|SmallOrderPrimes.admG1| asserts, under the quantified \path|GatTies|, \path|LatTied|, and \path|GIdx| conditions,

\[
 \begin{split}
 N&\subseteq p^{\max(0,\lceil\tau_P\rceil)}M_{\rm ram},\\
 \ell_O(M_{\rm ram}/N)&<\infty,\qquad\ell_O(M_{\rm ram}/N)\le e_p d_M(\tau_P+\operatorname{io}_P),\\
 f(N)&\subseteq
 \pi^{\left\lceil e_p\left(\sum_j\operatorname{gainP}_j+
       \sum_j(F_j-E_{u,j})\right)\right\rceil}B,
 \qquad d_M=\operatorname{Module.finrank}(O,M)=G.d.
 \end{split}\tag{2.3}
\]

The parameter $d_M$ is literally the source's natural-number
rank \texttt{G.d = Module.finrank G.O G.M}, defined in
\path|SharpRestriction/GrowthBounds/DataG.lean|, lines 188--194.
It occurs in the stated length bound without identifying it
with the rank of $M_{\rm ram}$. The later CM discriminant parameter
$d$ is a different quantity.

The last line is universally quantified over \textbf{every vector of \(N\)}. It is neither an assertion only about a determinant nor one about a single theta evaluation. It is also not an inclusion for every vector of the original \(M\).

The source's \path|UnitBoundLowTerm| is itself a uniform assertion for every \(F\in M\), after applying each designated product \(y_{V'}\). \path|TopChoiceBounds.bg_IM_of_Ulow| proves the resulting inclusion for the whole span \(I\): first it checks the displayed generators, then uses linearity for addition and scalar multiplication. \path|bg_NP_of_Ulow| restricts it to \(N\). In particular the alternative term

\[
 \operatorname{lowU}=-m\,v_p(D\Gamma)
\]

gives an inclusion on \(I\) whenever its \path|UnitBoundLowTerm| is assumed.

\textbf{Dependency scope.} \path|CornerGLContinuity.GrowthHypotheses| includes \path|ServiceLowerValuation|, which already states the required \(N\)-inclusion with the difference term in (2.3). Accordingly its \path|gluing_statement_holds| is not an independent derivation of that difference-term estimate from lattice volume alone. \path|FormalGroupSystem.LowTermBound| and \path|Conclusion5|, and the local-system comparison in \path|Joints| and \path|JointsTw|, also use universal vector bounds. Their geometric realization hypotheses still have to be verified for a transported system. These geometric realization hypotheses are not discharged by the algebraic transport.

\begingroup\small\raggedright
\textbf{Source locators:} \path|SharpRestriction/GrowthBounds/IsResidueClassZero.lean|, lines 411--428 and 459--472; \path|SharpRestriction/GrowthBounds/SmallOrderPrimes.lean|, lines 198--230; \path|SharpRestriction/GrowthBounds/CornerGLContinuity.lean|, lines 145--172 and 242--247; \path|SharpRestriction/GrowthBounds/TopChoiceBounds.lean|, lines 1464--1466 and 1586--1650. All source paths below are relative to the pinned repository's src/Formalism directory; the bibliographic entry gives its exact commit.\par\endgroup

\section{What a compatible averaging projector preserves}\label{what-a-compatible-averaging-projector-preserves}

\textbf{Proposition 3.1.} Let a finite group \(H\) act \(E\)-linearly on \(V_0,V_1\), with \(H\)-stable lattices \(N\subset V_0\) and \(B\subset V_1\). Let \(f:V_0\to V_1\) intertwine those actions. For \(a\in\mathbb Z\),

\[
 f(N)\subseteq\pi^a B
 \quad\Longrightarrow\quad
 f(e_HN)\subseteq\pi^a e_HB,\qquad
 e_H=|H|^{-1}\sum_{h\in H}h.\tag{3.1}
\]

\textbf{Proof.} \(fe_H=e_Hf\), so \(f(e_HN)=e_Hf(N)\subseteq e_H(\pi^aB)=\pi^ae_HB\). This is an exact inclusion with no exponent loss. If \(p\nmid|H|\), \(e_H\) is integral and \(e_HN=N^H\). At primes dividing \(|H|\), the same statement concerns the fractional lattices \(e_HN,e_HB\); replacing one of these by an integral invariant lattice is a separate operation and can introduce an index. \(\square\)

The index part of (2.3) needs separate attention even when \(p\nmid|H|\). A total normalized length bound need not have the same normalized bound on the invariant part. For example take \(M=O[H]\), \(|H|=r>1\) a unit in \(O\), and

\[
 N=\pi^a M^H\oplus\ker(e_H|_M),\qquad a>0.
\]

Then \(M=M^H\oplus\ker(e_H|_M)\), \(\ell(M/N)=a\), and \(\ell(M^H/N^H)=a\). Dividing by ranks gives \(a/r\) before projection and \(a\) after it. This is a concrete failure of a \textbf{normalized length inference}, not a failure of (3.1). It identifies a second quantity that a successful source-compatible projector must control.

\section{A coherent transported-copy construction and its exact output}\label{a-coherent-transported-copy-construction-and-its-exact-output}

Let \(H\) be any finite group. Start with finite-dimensional \(E\)-vector spaces \(V_j\), full lattices \(M_j\), a full sublattice \(N\subseteq M_0\) with finite-length quotient, and \(E\)-linear maps \(f_j:V_j\to V_{j+1}\). Assume \(M_0\) has positive rank when normalizing its index. Suppose for each \(h\in H\) an isomorphic transported copy \(V_{h,j}\) is chosen, with isomorphism \(T_{h,j}:V_j\to V_{h,j}\). Put

\[
 \widetilde V_j=\bigoplus_{h\in H}V_{h,j},\quad
 \widetilde M_j=\bigoplus_hT_{h,j}M_j,\quad
 \widetilde N=\bigoplus_hT_{h,0}N,
\]

and on the block indexed by \(h\) define

\[
 U_a=T_{ah,j}T_{h,j}^{-1},\qquad
 \widetilde f_j|_{V_{h,j}}=T_{h,j+1}f_jT_{h,j}^{-1}.\tag{4.1}
\]

\textbf{Proposition 4.1.} These are coherent group actions, all \(\widetilde f_j\) intertwine them, and every lattice inclusion for a composite of the \(f_j\) transfers with exactly the same exponent. The normalized input index is preserved both on the full direct sum and on its \(H\)-invariant projected lattice.

\textbf{Proof.} On the \(h\)-block, \(U_aU_b=T_{abh,j}T_{bh,j}^{-1}T_{bh,j}T_{h,j}^{-1}=U_{ab}\), and the analogous cancellation proves the intertwining equality. Define

\[
 \Phi_j:E[H]\otimes_E V_j\longrightarrow\widetilde V_j,\qquad
 [h]\otimes v\longmapsto T_{h,j}v.
\]

Then \(U_a\) becomes left translation on \(E[H]\) tensored with the identity, and \(\widetilde f_j\) becomes \(1\otimes f_j\). Lattices become \(O[H]\otimes M_j\), respectively \(O[H]\otimes N\). Inclusions are now componentwise. Full lengths and ranks both multiply by \(r=|H|\). The projected module is the diagonal copy generated by \(r^{-1}\sum_h[h]\); both its input and target lattices use this same scalar, so its quotient is isomorphic to \(M_0/N\), even when \(r\) is not a unit. Its rank is the original rank. This proves the index assertions as well. \(\square\)

The construction has no cocycle obstruction. In a polynomial tensor realization, a permutation \(h\) of the embedding labels gives the strict variable substitution \(X_{\tau,ij}\mapsto X_{h\tau,ij}\). Such substitutions compose exactly, carry tensor products to permuted tensor products, and carry intersections over all lattice frames to the corresponding permuted intersections. Transporting the whole datum therefore supplies examples of \(T_{h,j}\). This does \textbf{not} say that the source's already fixed operators on two independently chosen data are equal to the transported operators in (4.1).

The construction's limitation is explicit. The projector \(e_H\) recovers one copy of the original system:

\[
 (e_H\widetilde V_j,e_H\widetilde M_j,e_H\widetilde f_j)
 \simeq (V_j,M_j,f_j).\tag{4.2}
\]

Thus it proves equivariance by adding copies, but does not by itself produce the fixed cubic arithmetic quotient.

Here is the precise trace issue. Let \(A\) be a finite abelian deck group, \(H\subset A\), and put \(q=|A/H|\). On the artificially enlarged system take

\[
 \mathcal F_p=\chi(p)L_{a_p}\otimes T_p,
\]

where \(L_{a_p}\) is translation on \(E[A]\), \(\chi\) is a genuine cubic character, and \(T_p\) is an auxiliary invertible operator on \(V\). Then

\[
 \operatorname{tr}(e_H\mathcal F_p^\nu)
 =\chi(p)^\nu\,q\,\mathbf{1}_{a_p^\nu\in H}
       \operatorname{tr}(T_p^\nu).\tag{4.3}
\]

\textbf{Proof.} The invariant regular factor is \(E[A/H]\). Translation by \(a_p^\nu\) fixes every basis vector if its quotient class is the identity and fixes none otherwise. The trace of a tensor product is the product of the traces. \(\square\)

For the desired quotient the extra factor in (4.3) is absent. If it is required to equal \(1\) for every positive power of an invertible \(T_p\), then \(V\) must have dimension one and \(T_p=1\). Indeed the identity

\[
 \sum_{\nu\ge1}\operatorname{tr}(T_p^\nu)z^\nu
 =\sum_i\frac{\alpha_i z}{1-\alpha_i z}
 =\frac{z}{1-z}
\]

forces the only nonzero eigenvalue to be \(1\), with multiplicity one, by its poles and residues. For a quotient where only multiples of a fixed order \(d_p\) are visible, apply the same argument to \(T_p^{d_p}\); it forces dimension one and \(T_p^{d_p}=1\), while \(T_p\) itself may be a nontrivial \(d_p\)-th root of unity.

There is no claim that a multiplicity \(J_0^D\) could not imply RH for \(J_0\). Taking \(T_p=1_V\) would indeed introduce a harmless constant multiplicity \(D\) at the level of an already established Euler product. The missing implication is that the source's gain operators then control the Frobenius trace of this artificially added regular-character object. Formula (4.3) determines exactly what needs proving; a change of labels alone does not supply that comparison. A genuine cubic scalar twist preserves local unit-valued lattice inclusions, whereas a formal central third root of unity unconnected to \(\chi_{269}\) does not determine that character.

\section{The natural highest-weight orbit is covariant, not commuting}\label{the-natural-highest-weight-orbit-is-covariant-not-commuting}

Let \(E\) be algebraically closed of characteristic zero, \(B=\prod_{\tau\in I}\mathrm{GL}_k(E)\), and \(V_\lambda\) an irreducible representation with dominant highest-weight tuple \(\lambda=(\lambda_\tau)_\tau\). Let a permutation group \(H\) act on \(I\). Form the direct sum over the \textbf{distinct} tuples in the orbit,

\[
 V_{\mathcal O}=\bigoplus_{\mu\in H\lambda}V_\mu.
\]

The strict tensor permutations identify \(V_\mu\to V_{h\mu}\) and satisfy

\[
 U_h\rho(b)U_h^{-1}=\rho(hb),\tag{5.1}
\]

where \(hb\) permutes the factors of \(B\). This is a semidirect-product action, not a commuting \(H\times B\)-action.

\textbf{Proposition 5.1.} Every endomorphism commuting with both \(B\) and this transitive \(H\)-action is scalar. Consequently there is no nonzero proper idempotent with those two commutation properties.

\textbf{Proof.} The highest-weight classification \cite[Section~27.3]{etingof} identifies distinct tuples with nonisomorphic irreducible \(B\)-modules. Any \(B\)-intertwiner between different summands is zero: its kernel and image would otherwise give an isomorphism between irreducibles. An endomorphism of one summand is scalar: over an algebraically closed field it has an eigenvalue, whose eigenspace is a nonzero invariant subspace. Thus a \(B\)-commuting endomorphism is scalar on each summand. Commutation with the transitive permutation action equates all these scalars. A scalar idempotent is \(0\) or \(1\). \(\square\)

One can also see the failure directly. Normalize highest vectors so that tensor permutations send one to the next. Their orbit sum is \(H\)-fixed. A diagonal torus element acts on its components by the distinct characters \(\mu\); choosing an element on which two characters differ produces a vector which is not \(H\)-fixed. Hence the averaging projector does not commute with the coefficient torus for a nontrivial orbit.

Repeated isomorphic copies enlarge the commutant and evade this proposition, but then they return to the construction of Section 4. The proposition does not forbid a projector commuting only with a much smaller particular collection of theta-chain maps. Producing and controlling such maps would be an additional construction. The source's finite general-linear operators \path|eOp|, \path|eMt|, and the central finite-group idempotent matched to them are not already a field-deck action; substituting them for \(U_h\) does not prove (5.1) is commutation.

\section{The local parity and the dual-edge repair}\label{the-local-parity-and-the-dual-edge-repair}

The pinned \path|UnramifiedQuadraticSetting| uses an even cyclic period \(n'\ge2\) and defines

\[
 J(\lambda_1,\ldots,\lambda_k)=(-\lambda_k,\ldots,-\lambda_1),
 \qquad \widehat\lambda_s=J^{s\bmod2}\lambda_s.
\]

Here \(J^2=1\), and \(\widehat\lambda_s\) is its \path|lamAt|. Under the cyclic relabelling \(s\mapsto s+c\pmod {n'}\), strict variable permutation preserves the effective tensor weight precisely when

\[
 \lambda^{(c)}_{s+c}=J^{c\bmod2}\lambda_s.\tag{6.1}
\]

\textbf{Proof.} The target factor must have effective weight \(\widehat\lambda_s\), so \(J^{s+c}\lambda^{(c)}_{s+c}=J^s\lambda_s\). Multiplication by \(J^{s+c}\) proves (6.1). Reduction modulo the even integer \(n'\) leaves parity unchanged. This also proves coherence under addition of the shifts. \(\square\)

At a source top position, the first and last coordinates are \(a_s,-b_s\). An odd shift sends them to \(b_s,-a_s\), interchanging the two source windows. Under the literal full family assumptions these windows are disjoint. In fact its scale parameters satisfy

\[
 a=x\exp(x^\epsilon),\qquad b'=a/X',\qquad
 X'=\exp(x^{\epsilon/3})>e\quad(x>1,\epsilon>0).
\]

Writing \(D=2n\delta_1x^\epsilon\mathrm{La}/k\ge0\), the family gives \(b_s\le b_{\rm top}=b'(1+D)\) and \(a_s\ge a\). The formula for \(\rho_B\), nonnegative constants, the equality \(\mathrm{La}=\mathrm{LaG}\), and the family bound \(\rho_B\le x^{-2}\) give \(4D\le x^{-2}\). Therefore

\[
 b_s\le (1+x^{-2}/4)b'<(5/4)b'<a\le a_s.\tag{6.2}
\]

An odd pure relabelling accordingly leaves the admissible top windows whenever a top position is transported. If a particular orbit contains no top position, this window argument makes no assertion about that orbit.

A natural attempted repair is contragredient duality. It does not preserve the designated Hecke edge type. Let \(W\) be a DVR with uniformizer \(p\), residue cardinality \(Q\), and let \(L'\subset L\) be lattices of rank \(k\), with \(pL\subset L'\) and \(\ell_W(L/L')=r\). With respect to a basis,

\[
 L=W^k,\quad L'=pW^r\oplus W^{k-r},\quad
 L^*=W^k\subset (L')^*=p^{-1}W^r\oplus W^{k-r}.
\]

After rescaling to a lower edge,

\[
 p(L')^*=W^r\oplus pW^{k-r}\subset L^*,\qquad
 \ell_W\bigl(L^*/p(L')^*\bigr)=k-r.\tag{6.3}
\]

Thus lattice duality exchanges lower type-\(r\) edges with upper type-\(r\) edges, or, after the central rescaling, lower type-\(k-r\) edges. More explicitly, for functions on lattices let

\[
 T_r^-F(L)=\sum_{\substack{pL\subset L'\subset L\\\ell(L/L')=r}}F(L'),
 \quad T_r^+F(L)=\sum_{\substack{L\subset M\subset p^{-1}L\\\ell(M/L)=r}}F(M).
\]

Using the canonical double-dual identification, the involution \(DF(L)=F(L^*)\) satisfies \(DT_r^-D=T_r^+\), while

\[
 T_r^+=T_{k-r}^- C_{p^{-1}},\qquad C_{p^{-1}}F(L)=F(p^{-1}L).\tag{6.4}
\]

Both identities follow from the bijections just written. They can be distinguished on delta functions supported at a single edge lattice; a scalar normalization cannot change \(r\) to \(k-r\). The source fixes \(r=r_G\) or \(r=1\) through \path|LatTied|, whereas \(k\) grows with scale and \(r_G\) depends on the fixed scale exponent. Thus \(r\ne k-r\) for sufficiently large literal hosts. Its chain can change the ambient rank as well. Duality would require a new comparison for the complementary edge operators and their weight normalizations. The existing uniform image inclusion alone does not supply that comparison.

This is an exact obstruction to two proposed transports---pure odd relabelling and its direct dual-lattice repair---not an exclusion of all intertwining correspondences.

\begingroup\small\raggedright
\textbf{Source locators:}
\path|SharpRestriction/GrowthBounds/UnramifiedQuadraticSetting.lean|,
lines 12--35, 54--86 and 100--138;
\path|SharpRestriction/GrowthBounds/DataGChain.lean|,
lines 14--16 and 27--38;
\path|LocalRingCardinality/DataParams|\textsubscript{2}\path|.lean|,
lines 27--45 and 76--87;
\path|SharpRestriction/FamCore.lean|, lines 123--130;
\path|SharpRestriction/ParConsts.lean|, lines 31--48;
\path|SharpRestriction/ElevenLe.lean|, lines 68--75.\par\endgroup

\section{Why adjoining the source's CM quadratic base does not remove this obstruction}\label{why-adjoining-the-sources-cm-quadratic-base-does-not-remove-this-obstruction}

The source's Galois coordinates have \(G_0=(\mathbb Z^\times)^2\), complex conjugation \(c_0=(-1,-1)\), and total conjugation \((c_0,1)\). The quadratic CM base \(K_d=\mathbb Q(\sqrt{-d})\) is detected by the second \(G_0\)-coordinate; the analogous \(K_q\) is detected by the first. For an unramified prime of \path|inertClass|, under the source's explicit \path|FactClassFrob| comparison interface,

\[
 \operatorname{Frob}_p=(c_0,h).
\]

But the one-step translation in \path|GatTies.gat_step| is \textbf{Frobenius multiplied by complex conjugation}. Consequently it equals

\[
 \operatorname{Frob}_p\,c=((1,1),h).\tag{7.1}
\]

It fixes both \(K_d\) and \(K_q\). Whenever this element belongs to the proposed projector subgroup \(H=\operatorname{Gal}(F/F_0)\), it still belongs to

\[
 H'=\operatorname{Gal}(F/F_0K_d).
\]

Therefore adding \(K_d\) does not turn this one-step local rotation into an even shift. This is a statement about any prime with the indicated class and membership. It does not assert the existence of such a prime in the moving interval \(p\le2x\). More generally, a bounded-size quotient of a growing cyclic \(2\)-group has nontrivial kernel; adjoining a field on which the pure \(A\)-coordinate already acts trivially cannot remove that kernel. No effective prime-distribution conclusion is used here.

There is a second literal constraint on the suggested \emph{fixed} quadratic enlargement. \path|Host.qd_big| requires

\[
 d>4tr+1,\qquad r=k+m-1.
\]

The rank hypotheses give \(m\ge2, t\ge k+m+4, k\ge2n\). Hence

\[
 d>4k^2\ge16n^2.
\]

\path|ElevenLe.SizeN| gives \(n\ge x^{2+799\epsilon/100}\), so

\[
 \boxed{d>16x^{4+799\epsilon/50}.}\tag{7.2}
\]

In particular a single fixed \(d\equiv7\pmod {16}\) cannot satisfy these literal host hypotheses along an unbounded family. At each individual scale the compositum \(F_0K_d\) still has bounded degree if \(F_0\) does; (7.2) does not deny that. It says the added quadratic field must vary, and (7.1) shows that this variation does not resolve the specific local parity problem.

\begingroup\small\raggedright
\textbf{Source locators:} \path|LocalRingCardinality/DataGatTies.lean|, lines 10--31; \path|LocalRingCardinality/DataHost.lean|, lines 41--54 and 220--227; \path|SharpRestriction/RamificationFrobenius/TotientGaloisClasses/ConjEMemHostFIff.lean|, lines 29--31, 60--66, 200--210; \path|SharpRestriction/RamificationFrobenius/TotientGaloisClasses/ZsmulGeneratorEqZeroOfProjection.lean|, lines 102 and 120--122; \path|LocalRingCardinality/Params.lean|, lines 12--32; \path|SharpRestriction/ElevenLe.lean|, lines 11 and 24--25.\par\endgroup

\section{Resulting scope for the cubic RH question}\label{resulting-scope-for-the-cubic-rh-question}

The attempted method has been narrowed concretely. The source inclusion in (2.3) is strong enough to survive a compatible averaging projector, so determinant divisibility is not the correct general objection. Coherent transported copies exist and preserve that gain and its normalized lattice cost; their projected object and prime-power traces are exactly (4.2) and (4.3). Natural highest-weight orbits provide semidirect covariance rather than the needed commuting action. The local alternating weights and the complementary-edge calculation specify why two natural ways to repair this fail within the stated source data. Adding the source CM base does not repair the local step, and it cannot be held fixed under the literal size conditions.

A successful continuation would therefore require a theta-chain intertwiner, on the correctly projected selected lattices, whose arithmetic Frobenius traces give the fixed cubic quotient (or a controlled positive integral power of it), together with the projected normalized cost estimate. No prime-sum cancellation estimate or RH assertion follows merely from the transports proved here. These are proved algebraic conclusions and source-interface findings; no numerical evidence is being offered as proof.

\section*{Method and dedication}
The mathematics and text were produced by GPT-6 Astra in ChatGPT Pro
mode. A separate instance of the same model independently reconstructed
all six statement groups from their statements alone before reading the
author's proof. That reconstruction includes the elementary arguments
and the selected source interfaces at the pinned commit, with the
highest-weight classification as a named external input. Another
instance completed a targeted reading of the full text.
The source-specific conclusions retain their stated comparison and
geometric realization hypotheses. The pinned Lean source was inspected
but not compiled; no verification of every transitive dependency is claimed.
No human has checked this text.

The original contributions in this paper are dedicated to the public
domain, to the extent permitted by law, under
\href{https://creativecommons.org/publicdomain/zero/1.0/}{CC0 1.0}.
This dedication does not apply to cited third-party work.
The accompanying research context is the
\href{\workbenchurl}{zeta and L-functions workbench}, within the
\href{\recordurl}{mathematics workbenches collection}.

\begin{thebibliography}{9}
\bibitem{source}
CocoSgt (2026).
\emph{riemann-hypothesis-lean}, source code at commit
\path|7cbf7a0c32fca1badf05931879a37c3f8a2b03b4|.
\url{https://github.com/CocoSgt/riemann-hypothesis-lean/tree/7cbf7a0c32fca1badf05931879a37c3f8a2b03b4}.
Exact source-file and line locators are given in the text.
The source was inspected at the definition and selected-proof level;
it was not compiled, and no verification of every dependency is claimed.
\bibitem{etingof}
Etingof, P. (2024).
\emph{Lie groups and Lie algebras II: Lecture notes},
Section 27.3, pp.\ 146--147.
Massachusetts Institute of Technology.
\url{https://live.ocw.mit.edu/courses/18-755-lie-groups-and-lie-algebras-ii-spring-2024/mit18_755_s24_lec01.pdf}.
The highest-weight classification is the external input used
to distinguish the irreducible coefficient modules.
\end{thebibliography}
\end{document}
